\documentclass[12pt]{article}

\usepackage[T1]{fontenc}
\usepackage[utf8]{inputenc}
\usepackage{lmodern}
\usepackage[margin=1in]{geometry}
\usepackage{microtype}
\usepackage{amsmath,amssymb,amsthm,mathtools}
\usepackage{array,booktabs,enumitem}
\usepackage[colorlinks=true,linkcolor=blue,citecolor=blue,urlcolor=blue]{hyperref}
\usepackage[capitalize,nameinlink,noabbrev]{cleveref}

\newtheorem{theorem}{Theorem}[section]
\newtheorem{proposition}[theorem]{Proposition}
\newtheorem{lemma}[theorem]{Lemma}
\newtheorem{corollary}[theorem]{Corollary}
\newtheorem{condition}[theorem]{Condition}
\newtheorem{assumption}[theorem]{Assumption}
\theoremstyle{definition}
\newtheorem{definition}[theorem]{Definition}
\newtheorem{remark}[theorem]{Remark}
\newtheorem{example}[theorem]{Example}

\newcommand{\eps}{\varepsilon}
\newcommand{\R}{\mathbb{R}}
\newcommand{\N}{\mathbb{N}}
\newcommand{\E}{\mathbb{E}}
\newcommand{\Pb}{\mathbb{P}}
\newcommand{\1}{\mathbf{1}}
\newcommand{\cA}{\mathcal{A}}
\newcommand{\cB}{\mathcal{B}}
\newcommand{\cE}{\mathcal{E}}
\newcommand{\cF}{\mathcal{F}}
\newcommand{\cM}{\mathcal{M}}
\newcommand{\cN}{\mathcal{N}}
\newcommand{\cP}{\mathcal{P}}
\newcommand{\cQ}{\mathcal{Q}}
\newcommand{\cT}{\mathcal{T}}
\newcommand{\cS}{\mathcal{S}}
\newcommand{\cX}{\mathcal{X}}
\DeclareMathOperator{\conv}{conv}
\DeclareMathOperator{\supp}{supp}
\DeclareMathOperator{\spn}{sp}
\DeclareMathOperator{\Out}{Out}
\DeclareMathOperator{\BR}{BR}

\title{\textbf{Uniform Equilibrium from an Augmented Metastable Tree\\under Robust Nodewise Irreducibility}}
\author{Pedro Aldighieri}
\date{March 2026}

\begin{document}
\maketitle

\begin{abstract}
Neyman's problem asks whether every finite stochastic game admits a uniform
$\eps$-equilibrium.  The full conjecture remains open.  This paper closes the
\emph{downstream} part of the route based on a compact asymptotic tree object:
we work from a finite universal $\eta$-tree representation of the game and prove
that, once this representation is available, the route yields a public
finite-memory uniform $\eps$-equilibrium under a nodewise irreducibility
hypothesis.

The main repair is structural.  The v4 manuscript used a raw stationarity system
for node occupations.  Summing those equations forced zero exit mass and made the
feasible set empty at exit-capable nodes.  We replace that system by a
\emph{regenerated augmented system}: each node is completed by cemetery states
$\partial_b$ indexed by boundary states, and by a fixed bookkeeping regeneration
kernel that closes the local flow.  The correct augmented balance equations are
then consistent with strictly positive exit mass.  The feasible augmented set is
still the graph of the linear exit map,
\[
\cM_C^{\mathrm{aug}} = J_C(\cP_C),\qquad J_C(\mu)=(\mu,L_C\mu),
\]
but now the internal feasible set $\cP_C$ is nonempty at exit-capable nodes.

We also repair the local game and the global fixed point.  At each node and for a
fixed external environment, every player faces a finite average-reward MDP against
stationary mixed kernels of the other players; Bellman certificates therefore
exist.  A local stationary Nash equilibrium exists by Kakutani on the compact
product of statewise mixed-action simplices.  The global correspondence is defined
against \emph{frozen continuation values}: continuation values are read from the
reference package $x$, while the response package $y$ is allowed to differ from
$x$.  Self-consistency is imposed only at the fixed point, not beforehand.

Finally, we replace the spurious exponential local-controller estimate by the
polynomial rate $O(L^{-1/2})$, which is sufficient for the final
$\eps$-equilibrium argument.  We define the deviation MDP coefficients explicitly,
prove the coefficient-level Lipschitz estimate with the sharp factor
$1+\spn(h)$, and formulate condition~\textnormal{(a)} directly in the primitive
transition data.  Under that condition the bias span is uniformly bounded, and
the realized profile satisfies the regret estimate
\[
\eta + C_\eta\delta + \frac{2H_\eta}{M_*} + \frac{K_\eta}{T}.
\]
Choosing the parameters in that order yields the theorem.
\end{abstract}

\tableofcontents

\section{Introduction}

\subsection{Neyman's problem}

Shapley introduced stochastic games in 1953 and proved existence of stationary
equilibria for the discounted evaluation \cite{Shapley1953}.  The long-run average
evaluation is much subtler.  For zero-sum stochastic games, Mertens and Neyman
proved existence of the uniform value \cite{MertensNeyman1981}.  For two-player
non-zero-sum stochastic games, Vieille established the existence of uniform
$\eps$-equilibria \cite{Vieille2000I,Vieille2000II}.  Beyond two players the full
problem remains open; see Neyman's survey chapter, Solan's chapter on the
multi-player problem, Vieille's survey, and the handbook treatment
\cite{Neyman2003,Solan2003,Vieille2002,MertensSorinZamir2003}.  Positive results
are known in special subclasses, including three-player absorbing games
\cite{Solan1999}, correlated equilibrium constructions \cite{SolanVieille2002},
cyclic Markov equilibria phenomena \cite{FleschThuijsmanVrieze1997}, and a number
of extensions for nonzero-sum stochastic games in general state spaces
\cite{Nowak2003}.

The central question is:
\begin{quote}
\textbf{Neyman's open problem.}
For every finite stochastic game and every $\eps>0$, does there exist a strategy
profile that is an $\eps$-equilibrium simultaneously for all sufficiently large
finite horizons?
\end{quote}

This paper does not claim to settle the conjecture in full generality.  Its scope
is narrower and more precise: we start from the finite universal $\eta$-tree
representation produced by the route-D' reduction and show that, once that object
is available, the remainder of the route is correct under a robust nodewise
irreducibility condition.  The point of the rewrite is not to alter the theorem's
scope but to make every definition used downstream mathematically correct and every
stated proof complete.

\subsection{Main theorem}

We first record the standing scope.  The compact asymptotic tree representation is
an input to the present paper.

\begin{assumption}[Standing tree input]
\label{ass:tree-input}
For every $\eta\in(0,1)$, the game is endowed with a finite universal
$\eta$-tree representation as specified in \Cref{ass:node-data}.  All theorems in
Sections~\ref{sec:augmented}--\ref{sec:proof-main} are proved relative to that
finite-dimensional input.
\end{assumption}

Under this standing input the paper proves the following theorem.

\begin{theorem}[Main theorem]
\label{thm:main}
Let
\[
G=(N,S,(A_i)_{i\in N},u,P)
\]
be a finite stochastic game, and assume \Cref{ass:tree-input}.  Assume further
that condition~\textnormal{(a)} from \Cref{cond:a} holds for the node kernels of
the universal $\eta$-trees.  Then for every $\eps>0$ there exist a public
finite-memory strategy profile $\sigma^\eps$ and $T_0\in\N$ such that for every
$T\ge T_0$, every initial state $s\in S$, every player $i\in N$, and every
unilateral deviation $\tau_i$,
\[
\gamma_i^T(s,\sigma^\eps)
\ge
\gamma_i^T\bigl(s,(\tau_i,\sigma^\eps_{-i})\bigr)-\eps.
\]
In particular, every finite stochastic game equipped with a universal
$\eta$-tree representation satisfying condition~\textnormal{(a)} admits a public
finite-memory uniform $\eps$-equilibrium for every $\eps>0$.
\end{theorem}

\begin{remark}
The theorem is intentionally scoped to the downstream part of the compact-tree
route.  The full Neyman conjecture remains open because the route still requires,
in the background, construction of the universal $\eta$-tree representation for an
arbitrary game.  The purpose of the present rewrite is to show that the remaining
argument from that representation to a uniform equilibrium is mathematically
correct once the six foundational defects of v4 are repaired.
\end{remark}

\subsection{The six repairs}

The rewrite makes six structural changes.

\begin{enumerate}[label=\textnormal{(\arabic*)},leftmargin=2.4em]
\item The node occupation polytope is redefined on a regenerated augmented system,
so positive exit mass is compatible with the balance equations.
\item Local equilibrium existence is proved as a genuine multiplayer statement,
not inferred from one-player LP duality.
\item Continuation self-consistency is imposed only at the fixed point, not while
defining the response correspondence.
\item The local controller estimate is stated with the polynomial rate
$O(L^{-1/2})$, which is exactly what the final theorem needs.
\item The deviation MDP coefficients are defined explicitly, and the Bellman
transfer uses coefficient-level perturbation with the correct factor
$1+\spn(h)$.
\item Condition~\textnormal{(a)} is stated directly in the primitive transition
data and illustrated by an explicit example class.
\end{enumerate}

\subsection{Roadmap}

Section~\ref{sec:model} recalls the game model and states the finite-dimensional
node data used by the compact-tree route.  Section~\ref{sec:augmented} introduces
the regenerated augmented node system and proves the graph theorem for exit fluxes.
Section~\ref{sec:local} formulates the local average game and establishes local
best-response and local equilibrium existence.  Section~\ref{sec:global} defines
the global correspondence with frozen continuation values and proves a Kakutani
fixed point theorem.  Section~\ref{sec:realization} realizes the fixed point by
public finite-memory block controllers.  Section~\ref{sec:proof-main} proves the
bias-span estimate, the deviation-value Lipschitz lemma, the blockwise Bellman
bound, and finally \Cref{thm:main}.  The appendices record the finite-MDP facts,
a uniform block estimate for finite Markov chains, and the final aggregation
bookkeeping.

\section{Model and standing asymptotic input}
\label{sec:model}

\subsection{Finite stochastic games}

A finite stochastic game is a tuple
\[
G=(N,S,(A_i)_{i\in N},u,P)
\]
where $N=\{1,\dots,n\}$ is the finite set of players, $S$ is the finite state
space, $A_i(s)$ is the finite action set of player $i$ at state $s$, the stage
payoff of player $i$ is a function
\[
u_i:S\times \prod_{j\in N}A_j(s)\to [0,1],
\]
and the transition kernel is
\[
P(\cdot\mid s,a)\in \Delta(S)
\qquad (s\in S,\ a\in \prod_{j\in N}A_j(s)).
\]
The public history at stage $t$ is
\[
h_t=(s_1,a_1,s_2,a_2,\dots,s_t).
\]
A behavioral strategy of player $i$ assigns to every public history $h_t$ a
probability distribution on $A_i(s_t)$.

For a horizon $T\ge 1$, an initial state $s$, and a profile $\sigma$, the average
payoff of player $i$ is
\[
\gamma_i^T(s,\sigma)
=
\E_{s,\sigma}\Bigl[\frac1T\sum_{t=1}^T u_i(s_t,a_t)\Bigr].
\]

\begin{definition}
A profile $\sigma$ is an \emph{$\eps$-equilibrium for horizon $T$} if for every
player $i$, every initial state $s$, and every unilateral deviation $\tau_i$,
\[
\gamma_i^T(s,\sigma)
\ge
\gamma_i^T\bigl(s,(\tau_i,\sigma_{-i})\bigr)-\eps.
\]
It is a \emph{uniform $\eps$-equilibrium} if there exists $T_0$ such that the same
inequality holds for every $T\ge T_0$.
\end{definition}

\subsection{Node data from the tree representation}

The present paper starts after the route-D' reduction has produced a finite rooted
universal $\eta$-tree.  We record exactly the node data used later.

\begin{assumption}[Finite node data]
\label{ass:node-data}
Fix $\eta\in(0,1)$.  The universal $\eta$-tree $\cT_\eta$ is finite.  Every node
$C\in \cT_\eta$ carries the following data.
\begin{enumerate}[label=\textnormal{(D\arabic*)},leftmargin=2.8em]
\item A finite phase-state set
\[
\Xi_C=\bigsqcup_{k\in\mathbb Z/d_C\mathbb Z}\{k\}\times S_C^k.
\]
\item For each player $i$ and phase-state $\xi\in\Xi_C$, a finite action set
$A_i(\xi)$.
\item A finite boundary set $B_C$ and a successor map $b\mapsto C[b]$ sending
boundary states to child nodes or terminal labels.
\item Internal transition coefficients
\[
P_C^0(\xi'\mid \xi,a)\in[0,1],
\qquad \xi,\xi'\in \Xi_C,
\ a\in \prod_{j\in N}A_j(\xi),
\]
and exit coefficients
\[
\lambda_C(b\mid \xi,a)\in[0,1],
\qquad b\in B_C,
\]
such that
\[
\sum_{\xi'\in\Xi_C}P_C^0(\xi'\mid \xi,a)+\sum_{b\in B_C}\lambda_C(b\mid\xi,a)=1.
\]
\item A bookkeeping regeneration kernel
\[
\rho_C(\cdot\mid b)\in \Delta(\Xi_C)
\qquad (b\in B_C).
\]
\item A universal bias box constant $B_\eta^0<\infty$ such that every local Bellman
certificate used later can be normalized inside
$[-B_\eta^0,B_\eta^0]^{\Xi_C}$.
\end{enumerate}
\end{assumption}

\begin{remark}
The regeneration kernel $\rho_C$ is a bookkeeping device.  It is used only to
close the local balance equations on the augmented phase system.  Continuation
values are still carried by the cemetery labels $\partial_b$ associated with the
boundary states.  The physical role of $\rho_C$ is therefore normalization, not a
claim about what the original game does after leaving the node.
\end{remark}

\subsection{Pure and mixed stationary kernels inside a node}

Fix a node $C$.  A \emph{pure stationary selector} of player $i$ in node $C$ is a
map
\[
f_i: \Xi_C\to \bigsqcup_{\xi\in\Xi_C}A_i(\xi)
\quad\text{with }f_i(\xi)\in A_i(\xi).
\]
The finite set of such selectors is denoted by $F_{i,C}$.  A \emph{stationary
mixed kernel} of player $i$ is an element of
\[
\Sigma_{i,C}:=\prod_{\xi\in\Xi_C}\Delta(A_i(\xi)).
\]
We write $\Sigma_C:=\prod_{i\in N}\Sigma_{i,C}$.

\begin{remark}
A point of $\Sigma_{i,C}$ may be viewed either as a statewise mixed action kernel
or, equivalently, as a product distribution over the finite selector set
$F_{i,C}$.  We work with statewise mixed kernels because the deviation MDP
coefficients are affine in those kernels; this is the precise form needed later.
The local equilibrium existence result in Section~\ref{sec:local} is equivalent to
Nash's theorem for the finite selector game.
\end{remark}

\section{The regenerated augmented node system}
\label{sec:augmented}

\subsection{Augmented phase space and balance equations}

Fix a node $C\in\cT_\eta$.  Its augmented phase space is
\[
\Xi_C^{\mathrm{aug}}:=\Xi_C\sqcup \partial B_C,
\qquad
\partial B_C:=\{\partial_b:b\in B_C\}.
\]
The states $\partial_b$ are cemetery labels.  Inside the augmented bookkeeping
chain they are followed by the regeneration kernel $\rho_C(\cdot\mid b)$ from
\Cref{ass:node-data}.  This is the correction that prevents the balance equations
from forcing zero exit mass.

For a stationary mixed kernel $\sigma_C\in\Sigma_C$, define the induced internal
kernel
\[
Q_C^{\sigma}(\xi'\mid \xi)
:=
\sum_{a\in\prod_j A_j(\xi)} \sigma_C(a\mid \xi)
P_C^0(\xi'\mid \xi,a),
\]
and the induced exit kernel
\[
\Lambda_C^{\sigma}(b\mid\xi)
:=
\sum_{a\in\prod_j A_j(\xi)}\sigma_C(a\mid\xi)\lambda_C(b\mid\xi,a).
\]
The corresponding augmented kernel on $\Xi_C^{\mathrm{aug}}$ is
\begin{equation}
\label{eq:aug-kernel}
\widehat Q_C^{\sigma}(z\mid y)
=
\begin{cases}
Q_C^{\sigma}(\xi'\mid \xi), & y=\xi\in\Xi_C,\ z=\xi'\in\Xi_C,\\[0.3em]
\Lambda_C^{\sigma}(b\mid \xi), & y=\xi\in\Xi_C,\ z=\partial_b,\\[0.3em]
\rho_C(\xi'\mid b), & y=\partial_b,\ z=\xi'\in\Xi_C,\\[0.3em]
0, & y=\partial_b,\ z=\partial_{b'}.
\end{cases}
\end{equation}
This chain is finite.  For the fixed-point kernels used in the final proof,
condition~\textnormal{(a)} implies irreducibility of the internal chain and hence
uniqueness of the invariant distribution of the regenerated augmented chain.

If $\pi_C^{\sigma,\mathrm{aug}}$ denotes the invariant law of
$\widehat Q_C^{\sigma}$, define the internal occupation measure and exit mass by
\[
\mu_C^{\sigma}(\xi,a):=\pi_C^{\sigma,\mathrm{aug}}(\xi)\sigma_C(a\mid\xi),
\qquad
\beta_C^{\sigma}(b):=\pi_C^{\sigma,\mathrm{aug}}(\partial_b).
\]
The pair is denoted by
\[
\bar\mu_C^{\sigma}=(\mu_C^{\sigma},\beta_C^{\sigma}).
\]

\begin{definition}
\label{def:PC}
The \emph{internal feasible set} of node $C$ is the set of all
$\mu\in\R_+^{\Omega_C}$ for which there exists $\beta\in\R_+^{B_C}$ such that
\begin{align}
\sum_{(\xi,a)\in\Omega_C}\mu(\xi,a)+\sum_{b\in B_C}\beta(b)&=1,
\label{eq:norm-PC}
\\
\pi_\mu(\xi')
&=
\sum_{(\xi,a)\in\Omega_C}\mu(\xi,a)P_C^0(\xi'\mid\xi,a)
+
\sum_{b\in B_C}\beta(b)\rho_C(\xi'\mid b)
\qquad (\xi'\in\Xi_C),
\label{eq:internal-balance}
\\
\beta(b)
&=
\sum_{(\xi,a)\in\Omega_C}\mu(\xi,a)\lambda_C(b\mid\xi,a)
\qquad (b\in B_C),
\label{eq:beta-balance}
\end{align}
where
\[
\pi_\mu(\xi):=\sum_{a\in\prod_j A_j(\xi)}\mu(\xi,a).
\]
We denote this set by $\cP_C$.
\end{definition}

\begin{remark}
If one drops the regeneration term in \eqref{eq:internal-balance} and keeps only
raw internal transitions, summing the balance equations over $\xi'\in\Xi_C$
forces $\sum_b\beta(b)=0$.  That was the critical defect in v4.  The regenerated
augmented system is exactly what restores the missing mass.
\end{remark}

\subsection{The graph theorem for exit fluxes}

Define the linear exit map
\begin{equation}
\label{eq:LC-def}
L_C:\R^{\Omega_C}\to \R^{B_C},
\qquad
(L_C\mu)(b):=
\sum_{(\xi,a)\in\Omega_C}\mu(\xi,a)\lambda_C(b\mid\xi,a).
\end{equation}
The associated regeneration operator is
\begin{equation}
\label{eq:RC-def}
R_C:\R^{B_C}\to \R^{\Xi_C},
\qquad
(R_C\beta)(\xi'):=\sum_{b\in B_C}\beta(b)\rho_C(\xi'\mid b).
\end{equation}

\begin{proposition}
\label{prop:PC-polytope}
For every node $C$, the set $\cP_C$ is a nonempty compact convex polytope.
Moreover, $\mu\in\cP_C$ if and only if
\begin{align}
\sum_{(\xi,a)}\mu(\xi,a)+\sum_b (L_C\mu)(b)&=1,
\label{eq:PC-short1}
\\
\pi_\mu &= P_C^0\mu + R_C L_C\mu,
\label{eq:PC-short2}
\end{align}
where $(P_C^0\mu)(\xi'):=\sum_{(\xi,a)}\mu(\xi,a)P_C^0(\xi'\mid\xi,a)$.
\end{proposition}

\begin{proof}
The defining relations in \Cref{def:PC} are finitely many linear equalities and
inequalities in the variables $(\mu,\beta)$.  Hence the feasible set in
$\R_+^{\Omega_C}\times\R_+^{B_C}$ is a compact convex polyhedron; compactness
comes from \eqref{eq:norm-PC}.  Projecting onto the $\mu$-coordinates shows that
$\cP_C$ is a compact convex polytope.

If $\mu\in\cP_C$, the auxiliary vector $\beta$ is unique by
\eqref{eq:beta-balance}, namely $\beta=L_C\mu$.  Substituting this into
\eqref{eq:norm-PC} and \eqref{eq:internal-balance} gives
\eqref{eq:PC-short1}--\eqref{eq:PC-short2}.  Conversely, if
\eqref{eq:PC-short1}--\eqref{eq:PC-short2} hold and one sets $\beta=L_C\mu$, then
all three relations from \Cref{def:PC} are satisfied.
\end{proof}

\begin{definition}
\label{def:MCaug}
The \emph{augmented feasible set} of node $C$ is
\[
\cM_C^{\mathrm{aug}}
:=
\{(\mu,\beta)\in\cP_C\times\R_+^{B_C}:\beta=L_C\mu\}.
\]
The affine graph map is
\[
J_C:\cP_C\to \cM_C^{\mathrm{aug}},
\qquad J_C(\mu)=(\mu,L_C\mu).
\]
\end{definition}

\begin{theorem}[Graph theorem for the augmented node]
\label{thm:graph-theorem}
For every node $C$,
\[
\cM_C^{\mathrm{aug}} = J_C(\cP_C).
\]
The map $J_C$ is affine and injective.  In particular the exit coordinates are not
independent variables: once the internal part $\mu$ is fixed, the exit mass is
uniquely determined as $L_C\mu$.
\end{theorem}

\begin{proof}
The identity is immediate from \Cref{def:MCaug}.  Affineness follows from the
linearity of $L_C$, and injectivity is obvious from the first coordinate.
\end{proof}

\begin{corollary}
\label{cor:no-vp3}
There is no separate compatibility condition between internal occupation and exit
mass.  The feasible augmented set is exactly the graph of the linear exit map.
\end{corollary}

\begin{proof}
This is a restatement of \Cref{thm:graph-theorem}.
\end{proof}

\subsection{Actual local packages induced by stationary mixed kernels}

\begin{definition}
\label{def:actual-local-package}
For a stationary mixed kernel $\sigma_C\in\Sigma_C$, its \emph{actual augmented
occupation} is the pair $\bar\mu_C^{\sigma}$ induced by the invariant law of the
augmented kernel \eqref{eq:aug-kernel}.  By construction,
$\bar\mu_C^{\sigma}\in \cM_C^{\mathrm{aug}}$.
\end{definition}

\begin{proof}
Let $\pi=\pi_C^{\sigma,\mathrm{aug}}$ be the invariant law of
$\widehat Q_C^{\sigma}$.  Writing stationarity at internal states and at cemetery
states gives exactly \eqref{eq:norm-PC}--\eqref{eq:beta-balance} with
$\mu=\mu_C^{\sigma}$ and $\beta=\beta_C^{\sigma}$.
\end{proof}

\section{Local average-reward games and Bellman certificates}
\label{sec:local}

\subsection{Local continuation values}

Fix $\eta\in(0,1)$ and a node $C$.  Let $x$ be a package of continuation data on the
whole tree.  For each boundary state $b\in B_C$ the value $v_C^x(b)\in[0,1]^N$ is
read from the continuation coordinates of $x$: if $C[b]$ is a child node $D$, then
$v_C^x(b)$ is the continuation vector stored at $b$ in $x$ and required later to be
$\eta$-close to $g_D$; if $C[b]$ is a terminal label, then $v_C^x(b)$ is the
terminal continuation value attached to that label.  Throughout this section $x$ is
considered fixed.

\subsection{Deviation MDP coefficients}

This subsection repairs the coefficient map that was left implicit in v4.

Fix a node $C$, a player $i$, and a stationary mixed kernel
$\sigma_{-i}\in \prod_{j\ne i}\Sigma_{j,C}$.  The unilateral deviation MDP of
player $i$ has state space $\Xi_C$, action sets $A_i(\xi)$, and coefficients
\begin{align}
\label{eq:r-sigma}
r_{i,C}^{\sigma_{-i},x}(\xi,a_i)
&:=
\sum_{a_{-i}}\sigma_{-i}(a_{-i}\mid\xi)
\,u_i\bigl(\xi,(a_i,a_{-i})\bigr),
\\
\label{eq:Q-sigma}
Q_{i,C}^{\sigma_{-i}}(\xi'\mid \xi,a_i)
&:=
\sum_{a_{-i}}\sigma_{-i}(a_{-i}\mid\xi)
\,P_C^0\bigl(\xi'\mid\xi,(a_i,a_{-i})\bigr),
\\
\label{eq:q-exit-sigma}
\Lambda_{i,C}^{\sigma_{-i}}(b\mid\xi,a_i)
&:=
\sum_{a_{-i}}\sigma_{-i}(a_{-i}\mid\xi)
\,\lambda_C\bigl(b\mid\xi,(a_i,a_{-i})\bigr).
\end{align}
The continuation-adjusted Bellman operator is
\begin{equation}
\label{eq:Bellman-operator}
(T_{i,C}^{\sigma_{-i},x}h)(\xi)
:=
\max_{a_i\in A_i(\xi)}\Bigl[
 r_{i,C}^{\sigma_{-i},x}(\xi,a_i)
 +\sum_{\xi'\in\Xi_C}Q_{i,C}^{\sigma_{-i}}(\xi'\mid\xi,a_i)h(\xi')
 +\sum_{b\in B_C}\Lambda_{i,C}^{\sigma_{-i}}(b\mid\xi,a_i)v_{i,C}^x(b)
 \Bigr].
\end{equation}

\begin{lemma}
\label{lem:coeff-affine}
For fixed $C$, $i$, and $x$, the maps
\[
\sigma_{-i}\mapsto r_{i,C}^{\sigma_{-i},x},
\qquad
\sigma_{-i}\mapsto Q_{i,C}^{\sigma_{-i}},
\qquad
\sigma_{-i}\mapsto \Lambda_{i,C}^{\sigma_{-i}}
\]
are affine.  Consequently there exists a finite constant $K_{i,C}$ such that for
all $\sigma_{-i},\tau_{-i}$,
\begin{align}
\|r_{i,C}^{\tau_{-i},x}-r_{i,C}^{\sigma_{-i},x}\|_\infty
&\le K_{i,C}\|\tau_{-i}-\sigma_{-i}\|_1,
\label{eq:r-lip}
\\
\sup_{\xi,a_i}\|Q_{i,C}^{\tau_{-i}}(\cdot\mid\xi,a_i)-Q_{i,C}^{\sigma_{-i}}(\cdot\mid\xi,a_i)\|_{\mathrm{TV}}
&\le K_{i,C}\|\tau_{-i}-\sigma_{-i}\|_1,
\label{eq:Q-lip}
\\
\sup_{\xi,a_i}\|\Lambda_{i,C}^{\tau_{-i}}(\cdot\mid\xi,a_i)-\Lambda_{i,C}^{\sigma_{-i}}(\cdot\mid\xi,a_i)\|_{1}
&\le K_{i,C}\|\tau_{-i}-\sigma_{-i}\|_1.
\label{eq:Lambda-lip}
\end{align}
\end{lemma}

\begin{proof}
Each coefficient is a finite linear combination of the components of
$\sigma_{-i}(\cdot\mid\xi)$, with coefficients given by the primitive reward and
transition data.  This proves affineness.  The Lipschitz bounds follow because all
underlying arrays are bounded and finite.
\end{proof}

\subsection{Local best responses}

We now formulate the local best-response correspondence.  Since the local state and
action sets are finite, every unilateral local problem is a finite average-reward
MDP.

\begin{definition}
\label{def:local-BR-set}
For fixed $x$, node $C$, and player $i$, let $\mathfrak B_{i,C}(x)$ be the set of
triples $(\sigma_{i,C},g_{i,C},h_{i,C})$ such that
\begin{enumerate}[label=\textnormal{(BR\arabic*)},leftmargin=2.9em]
\item $\sigma_{i,C}\in \Sigma_{i,C}$;
\item $g_{i,C}\in[0,1]$ and $h_{i,C}\in[-B_\eta^0,B_\eta^0]^{\Xi_C}$;
\item $h_{i,C}$ is normalized by
\[
\min_{\xi\in\Xi_C}h_{i,C}(\xi)=0;
\]
\item for every $\xi\in\Xi_C$,
\begin{equation}
\label{eq:Bellman-ineq-def}
g_{i,C}+h_{i,C}(\xi)
\ge
r_{i,C}^{\sigma_{-i,C}^x,x}(\xi,a_i)
+
\sum_{\xi'}Q_{i,C}^{\sigma_{-i,C}^x}(\xi'\mid\xi,a_i)h_{i,C}(\xi')
+
\sum_b \Lambda_{i,C}^{\sigma_{-i,C}^x}(b\mid\xi,a_i)v_{i,C}^x(b)
\end{equation}
for every $a_i\in A_i(\xi)$;
\item equality holds in \eqref{eq:Bellman-ineq-def} for every action $a_i$ in the
support of $\sigma_{i,C}(\cdot\mid\xi)$.
\end{enumerate}
\end{definition}

\begin{proposition}
\label{prop:local-BR}
Assume condition~\textnormal{(a)}.  For every fixed $x$, node $C$, and player $i$, the set $\mathfrak B_{i,C}(x)$ is
nonempty, compact, and convex.  The correspondence
$x\mapsto \mathfrak B_{i,C}(x)$ has closed graph.
\end{proposition}

\begin{proof}
Fix $x$, $C$, and $i$.  The Bellman system
\eqref{eq:Bellman-ineq-def} is exactly the dual optimality system for the finite
average-reward MDP faced by player $i$ against the frozen environment encoded by
$x$.  By the finite-state average-reward LP duality recorded in
\Cref{prop:mdp-duality}, there exists an optimal stationary mixed kernel together
with a gain-bias pair satisfying \eqref{eq:Bellman-ineq-def} and the support
equalities with equality on the support.  Hence $\mathfrak B_{i,C}(x)$ is
nonempty.

Compactness follows because $\Sigma_{i,C}$ is a product of simplices and the gain
and bias coordinates are confined to compact boxes.  The defining equalities and
inequalities are closed.  For convexity, note that the local deviation MDP is
communicating on $\Xi_C$ under the standing node construction, so the optimal gain is
unique and the normalized optimal bias is unique.  Therefore, once $x$ is fixed,
there is a unique normalized optimal bias $h^*$ and a unique optimal gain $g^*$.  A
stationary mixed kernel is optimal if and only if at every state it assigns mass
only to actions attaining equality in the Bellman system for $(g^*,h^*)$.  The set
of such kernels is a product of simplices and hence convex.  Thus
$\mathfrak B_{i,C}(x)$ is convex.

For the graph property, let $x^m\to x$ and let
$(\sigma_i^m,g_i^m,h_i^m)\in \mathfrak B_{i,C}(x^m)$ converge to
$(\sigma_i,g_i,h_i)$.  The coefficient arrays from
\Cref{lem:coeff-affine} depend continuously on $x$, hence passing to the limit in
the Bellman equalities and inequalities shows
$(\sigma_i,g_i,h_i)\in\mathfrak B_{i,C}(x)$.
\end{proof}

\subsection{Local equilibrium existence}

\begin{proposition}[Local stationary Nash equilibrium]
\label{prop:local-equilibrium}
Assume condition~\textnormal{(a)}.  Fix $x$ and a node $C$.  There exists a stationary mixed kernel
$\sigma_C^*\in\Sigma_C$ such that for every player $i$,
$(\sigma_{i,C}^*,g_{i,C}^*,h_{i,C}^*)\in \mathfrak B_{i,C}(x)$ for some
$(g_{i,C}^*,h_{i,C}^*)$.  Equivalently, the local average game at node $C$ with
continuation values frozen at $x$ has a stationary Nash equilibrium, and every
player admits a Bellman certificate at that equilibrium.
\end{proposition}

\begin{proof}
Fix the continuation vector $v_C^x(\cdot)$ extracted from $x$.  For each player
$i$ and every opponents' kernel $\sigma_{-i,C}\in\prod_{j\ne i}\Sigma_{j,C}$, let
\[
\BR_{i,C}^x(\sigma_{-i,C})
:=
\{\tau_{i,C}\in\Sigma_{i,C}: \exists g,h\text{ with }
(\tau_{i,C},g,h)\text{ satisfying \eqref{eq:Bellman-ineq-def} against }
\sigma_{-i,C}\text{ and }v_C^x\}.
\]
By the same LP-duality argument as in \Cref{prop:local-BR}, each value
$\BR_{i,C}^x(\sigma_{-i,C})$ is nonempty, compact, and convex, and the graph of the
playerwise best-response correspondence is closed.  Define the product
correspondence on $\Sigma_C$ by
\[
\BR_C^x(\sigma_C):=\prod_{i\in N}\BR_{i,C}^x(\sigma_{-i,C}).
\]
Its values are therefore nonempty, compact, and convex, and its graph is closed.
Since $\Sigma_C$ is a compact convex product of simplices, Kakutani's theorem
\cite{Kakutani1941} yields a fixed point $\sigma_C^*\in\BR_C^x(\sigma_C^*)$.  For
each player $i$, choose a Bellman certificate $(g_{i,C}^*,h_{i,C}^*)$ associated
with the best-response relation at that fixed point.  This is exactly the required
local stationary Nash package.
\end{proof}

\begin{remark}
This proposition is the correct multiplayer replacement for the v4 appendix, which
contained only one-player LP duality.  One-player LP duality is still needed, but
only to identify each player's best-response set.  Existence of a simultaneous
local Nash package requires the fixed-point step above.
\end{remark}

\section{The global correspondence with frozen continuation values}
\label{sec:global}

\subsection{Ambient package space}

A global package contains, at each node and for each player, a stationary mixed
kernel and a Bellman certificate, together with nodewise continuation values.
There is \emph{no} self-consistency requirement in the ambient space.

\begin{definition}
\label{def:Aeta}
For fixed $\eta\in(0,1)$, the ambient package space $\cA_\eta$ consists of all
families
\[
x=\Bigl((\sigma_{i,C},g_{i,C},h_{i,C})_{i\in N},\ (v_C(b))_{b\in B_C}\Bigr)_{C\in\cT_\eta}
\]
such that:
\begin{enumerate}[label=\textnormal{(A\arabic*)},leftmargin=2.8em]
\item $\sigma_{i,C}\in \Sigma_{i,C}$ for all $i,C$;
\item $g_{i,C}\in[0,1]$ for all $i,C$;
\item $h_{i,C}\in[-B_\eta^0,B_\eta^0]^{\Xi_C}$ and
$\min_{\xi}h_{i,C}(\xi)=0$ for all $i,C$;
\item $v_C(b)\in[0,1]^N$ for all $C$ and $b\in B_C$.
\end{enumerate}
\end{definition}

\begin{proposition}
\label{prop:Aeta-compact}
For every $\eta\in(0,1)$, the set $\cA_\eta$ is compact and convex.
\end{proposition}

\begin{proof}
Each $\Sigma_{i,C}$ is a product of simplices, hence compact and convex.  The gain
coordinates lie in $[0,1]$, the bias coordinates lie in a compact box intersected
with a linear normalization hyperplane, and the continuation values lie in cubes.
Because the tree is finite, $\cA_\eta$ is a finite product of compact convex sets.
\end{proof}

\subsection{Global best response}

The key point is that continuation values are frozen from the reference package
$x$.  The response package $y$ is allowed to differ from $x$.

\begin{definition}
\label{def:global-BR}
For $x\in\cA_\eta$, define $\BR_\eta(x)$ to be the set of all packages
$y\in\cA_\eta$ such that:
\begin{enumerate}[label=\textnormal{(G\arabic*)},leftmargin=2.8em]
\item for every node $C$ and every player $i$,
\[
(\sigma_{i,C}^y,g_{i,C}^y,h_{i,C}^y)\in \mathfrak B_{i,C}(x);
\]
\item whenever $b\in B_C$ leads to a child node $D=C[b]$, one has
\begin{equation}
\label{eq:frozen-continuation}
\|v_C^y(b)-g_D^x\|_\infty\le \eta,
\end{equation}
where $g_D^x=(g_{i,D}^x)_{i\in N}$;
\item if $b\in B_C$ leads to a terminal label with continuation vector
$w(b)\in[0,1]^N$, then
\[
v_C^y(b)=w(b).
\]
\end{enumerate}
\end{definition}

\begin{remark}
The response package $y$ is assembled against \emph{frozen} continuation values and
frozen opponents from $x$.  There is no circularity here.  One does not ask $y$ to
be self-consistent while defining $\BR_\eta(x)$; that happens only at a fixed
point.  This is exactly the continuation repair missing in v4.
\end{remark}

\begin{theorem}[Global Kakutani fixed point]
\label{thm:kakutani-global}
Assume condition~\textnormal{(a)}.  For every $\eta\in(0,1)$, the correspondence
$\BR_\eta:\cA_\eta\rightrightarrows\cA_\eta$ has nonempty compact convex values and
closed graph.  Consequently it has a fixed point $x^*\in\cA_\eta$.
\end{theorem}

\begin{proof}
By \Cref{prop:Aeta-compact}, the domain $\cA_\eta$ is compact and convex.

Fix $x\in\cA_\eta$.  For every node $C$ and player $i$, the set
$\mathfrak B_{i,C}(x)$ is nonempty, compact, and convex by
\Cref{prop:local-BR}.  The continuation constraints
\eqref{eq:frozen-continuation} are closed convex constraints, and terminal
continuation coordinates are fixed singleton values.  Therefore $\BR_\eta(x)$ is a
finite product of nonempty compact convex sets, hence itself nonempty, compact, and
convex.

For the graph property, let $x^m\to x$ and $y^m\to y$ with
$y^m\in\BR_\eta(x^m)$ for all $m$.  For every node $C$ and player $i$ we have
$(\sigma_{i,C}^{y^m},g_{i,C}^{y^m},h_{i,C}^{y^m})\in\mathfrak B_{i,C}(x^m)$.  By
the closed-graph property from \Cref{prop:local-BR}, the limit belongs to
$\mathfrak B_{i,C}(x)$.  Passing to the limit in the continuation inequalities gives
\eqref{eq:frozen-continuation} for $y$ against $x$.  Hence $y\in\BR_\eta(x)$, so
$\BR_\eta$ has closed graph.

Closed graph plus compact values on a compact domain imply upper
hemicontinuity.  Kakutani's theorem \cite{Kakutani1941} yields a fixed point.
\end{proof}

\begin{corollary}
\label{cor:local-nash-at-fixed-point}
If $x^*\in\BR_\eta(x^*)$, then for every node $C$ the stationary mixed kernel
$\sigma_C^*:=(\sigma_{i,C}^*)_{i\in N}$ is a local stationary Nash equilibrium of
the node game with continuation values $v_C^{x^*}(\cdot)$, and each player has a
Bellman certificate $(g_{i,C}^*,h_{i,C}^*)$.
\end{corollary}

\begin{proof}
At a fixed point,
$(\sigma_{i,C}^*,g_{i,C}^*,h_{i,C}^*)\in\mathfrak B_{i,C}(x^*)$ for every $i,C$.
Thus each player's stationary mixed kernel is a best response to the others inside
node $C$, against the continuation values carried by $x^*$.  That is exactly the
local stationary Nash property.
\end{proof}

\section{Local controllers and global realization}
\label{sec:realization}

\subsection{The node controller}

Fix a fixed point $x^*\in\BR_\eta(x^*)$.  For each node $C$, write
\[
\sigma_C^*:=(\sigma_{i,C}^*)_{i\in N},
\qquad
\bar\mu_C^*:=(\mu_C^{\sigma_C^*},\beta_C^{\sigma_C^*}).
\]
The local controller for node $C$ is the public finite-memory device that, during a
block of length $L$, plays the stationary mixed kernel $\sigma_C^*$ on the phase
states of node $C$.

Let $\widehat\mu_{C,L}^{\mathrm{aug}}$ denote the empirical augmented occupation
vector over an $L$-block, normalized by block length.  The next theorem is the
repaired version of the local controller lemma: the relevant estimate is
$O(L^{-1/2})$, not a fake exponential remainder.

\begin{theorem}[Local controller lemma]
\label{thm:local-controller}
Assume condition~\textnormal{(a)}.  Fix a node $C$ and the fixed-point kernel $\sigma_C^*$.  There exists a constant
$K_C<\infty$ such that for every block length $L\ge 1$, every entry phase-state,
and every player $i$,
\begin{equation}
\label{eq:local-controller-occ}
\E\bigl[\|\widehat\mu_{C,L}^{\mathrm{aug}}-\bar\mu_C^*\|_1\bigr]
\le \frac{K_C}{\sqrt L},
\end{equation}
and
\begin{equation}
\label{eq:local-controller-payoff}
\left|
\E\Bigl[\frac1L\sum_{t=1}^L u_i^t\Bigr]-g_{i,C}^*
\right|
\le \frac{K_C}{\sqrt L}.
\end{equation}
Consequently, for every $\delta>0$ there exists $L_C(\delta)$ such that the block
controller is $\delta$-accurate in the sense of
\eqref{eq:local-controller-occ}--\eqref{eq:local-controller-payoff} whenever
$L\ge L_C(\delta)$.
\end{theorem}

\begin{proof}
Fix $C$.  Under condition~\textnormal{(a)}, the internal chain induced by the fixed
kernel $\sigma_C^*$ is irreducible, hence the regenerated augmented chain on
$\Xi_C^{\mathrm{aug}}$ is irreducible as well and has a unique invariant
probability $\pi_C^{\sigma_C^*,\mathrm{aug}}$.  Each coordinate of the empirical
occupation vector is the empirical average of a bounded function on that chain.
The uniform $L^1$ bound \eqref{eq:local-controller-occ} follows from the
finite-state block estimate in \Cref{prop:block-estimate}; because there are
finitely many coordinates, one may sum the coordinatewise bounds and absorb the
resulting constant into $K_C$.

To identify the target payoff, integrate the Bellman equalities at actions in the
support of $\sigma_{i,C}^*(\cdot\mid\xi)$ against the invariant occupation
$\bar\mu_C^*$.  The internal transition terms telescope under stationarity and the
cemetery terms contribute exactly the continuation coordinates carried by
$\beta_C^{\sigma_C^*}$.  One obtains
\[
g_{i,C}^*
=
\sum_{(\xi,a)}\mu_C^{\sigma_C^*}(\xi,a)u_i(\xi,a)
+
\sum_{b\in B_C}\beta_C^{\sigma_C^*}(b)v_{i,C}^{x^*}(b).
\]
Therefore the target value of the block-payoff functional at $\bar\mu_C^*$ is
precisely $g_{i,C}^*$.  Since the block-average payoff is a bounded linear
functional of the empirical augmented occupation vector, the occupation estimate
\eqref{eq:local-controller-occ} implies \eqref{eq:local-controller-payoff} after
increasing $K_C$ if necessary.
\end{proof}

\subsection{Frozen superexponential schedule}

\begin{definition}
A sequence $(M_r)_{r\ge 1}$ of positive integers is a
\emph{superexponential schedule} if $M_{r+1}\ge M_r^2$ for all $r$.  Its
\emph{freezing at level $R$} is the sequence $(M_r^{(R)})_{r\ge 1}$ defined by
\[
M_r^{(R)}=
\begin{cases}
M_r, & r<R,\\
M_R, & r\ge R.
\end{cases}
\]
We write $M_*:=M_R$ for the tail block length.
\end{definition}

\begin{lemma}
\label{lem:frozen-memory}
A frozen superexponential schedule uses only finite public memory.
\end{lemma}

\begin{proof}
Only the finitely many lengths $M_1,\dots,M_R$ appear.  To implement the schedule
one stores the current node, the current position inside the current block, the
current local memory state of the node controller, and the block index truncated at
$R$.  All of these coordinates range over finite sets.
\end{proof}

\subsection{Public routing and realization}

\begin{proposition}[Public routing]
\label{prop:public-routing}
Fix a fixed point $x^*$ and a frozen superexponential schedule.  There exists a
public finite-memory profile $\sigma_{\eta,\delta,M_*}$ with the following
properties.
\begin{enumerate}[label=\textnormal{(R\arabic*)},leftmargin=2.8em]
\item At each complete block played in node $C$, the profile runs the node
controller from \Cref{thm:local-controller} with block length equal to the current
schedule value.
\item The public memory stores the current node, the current block index, and the
current local controller state.
\item If a block emits a boundary state $b\in B_C$, the next block starts in the
child node $C[b]$; if no exit is emitted, the next block remains in node $C$.
\item The routing update is publicly observable and deterministic conditional on the
public memory and the observed block output.
\end{enumerate}
\end{proposition}

\begin{proof}
The node controller is public by construction.  The current node is public by
induction on completed blocks.  The block index evolves deterministically.  The
block output is part of the public phase history.  The successor map
$b\mapsto C[b]$ belongs to the node data, so the routing update is deterministic
once the output is observed.
\end{proof}

\begin{theorem}[Realization theorem]
\label{thm:realization}
Assume condition~\textnormal{(a)}.  Fix $\eta\in(0,1)$ and a fixed point $x^*\in\BR_\eta(x^*)$.  Let
$\sigma_{\eta,\delta,M_*}$ be the public profile obtained from
\Cref{prop:public-routing} by freezing a superexponential schedule at level $R$.
Then there exist constants $C_\eta,K_\eta<\infty$ such that:
\begin{enumerate}[label=\textnormal{(Q\arabic*)},leftmargin=2.8em]
\item if a complete block $B$ is played in node $C$ and has length $|B|\ge M_*$,
then
\begin{equation}
\label{eq:block-occupation-realization}
\E\bigl[\|\widehat\mu_B^{\mathrm{aug}}-\bar\mu_C^*\|_1\mid \cF_{\mathrm{in}(B)}\bigr]
\le \delta;
\end{equation}
\item for every player $i$ and every complete block $B$ played in node $C$,
\begin{equation}
\label{eq:block-payoff-realization}
\left|
\E\Bigl[\frac1{|B|}\sum_{t\in B}u_i^t\,
\Bigm|\,\cF_{\mathrm{in}(B)}\Bigr]-g_{i,C}^*
\right|
\le C_\eta\delta;
\end{equation}
\item for every horizon $T$,
\begin{equation}
\label{eq:realization-tail-bound}
\left|
\gamma_i^T(s,\sigma_{\eta,\delta,M_*})
-
\frac1T\sum_{B\subseteq [1,T]} |B|\,g_{i,C(B)}^*
\right|
\le C_\eta\delta + \frac{K_\eta}{T},
\end{equation}
where the sum is over complete blocks contained in $[1,T]$ and $C(B)$ denotes the
node of block $B$.
\end{enumerate}
\end{theorem}

\begin{proof}
For complete blocks of tail length at least $M_*$, the occupation estimate from
\Cref{thm:local-controller} and the choice of $M_*$ imply
\eqref{eq:block-occupation-realization}.  Since stage payoffs and continuation
terms are bounded linear functionals of the augmented occupation vector,
\eqref{eq:block-payoff-realization} follows after multiplying the occupation error
by a uniform Lipschitz constant depending only on the finite tree object.

To prove \eqref{eq:realization-tail-bound}, decompose the horizon into the complete
blocks contained in $[1,T]$ plus at most one initial incomplete fragment and one
final incomplete fragment.  The incomplete fragments contribute at most a constant
$K_\eta$ because the frozen prelude is finite and because stage payoffs are bounded
by $1$.  Summing \eqref{eq:block-payoff-realization} over the complete blocks and
dividing by $T$ yields \eqref{eq:realization-tail-bound}.
\end{proof}

\section{Condition \texorpdfstring{$(a)$}{(a)}, bias spans, and the proof of the main theorem}
\label{sec:proof-main}

\subsection{Condition \texorpdfstring{$(a)$}{(a)} in primitive data}

We now formulate the nodewise irreducibility hypothesis directly from the primitive
transition data.

\begin{condition}[Robust nodewise irreducibility]
\label{cond:a}
For every $\eta\in(0,1)$, every node $C\in\cT_\eta$, every player $i\in N$, every
statewise mixed kernel of the other players $\sigma_{-i}\in\prod_{j\ne i}\Sigma_{j,C}$,
and every pure action selection $a_i(\cdot)$ of player $i$, the internal kernel
\[
\xi'\mapsto \sum_{a_{-i}}\sigma_{-i}(a_{-i}\mid\xi)
P_C^0\bigl(\xi'\mid\xi,(a_i(\xi),a_{-i})\bigr)
\]
is irreducible on $\Xi_C$.
\end{condition}

\begin{remark}
Condition~\textnormal{(a)} has a direct operational meaning: no single player can,
by choosing actions state by state while the others keep their local kernels,
create a noncommunicating class inside a metastable node.
\end{remark}

\begin{example}
\label{ex:cond-a-trivial}
Suppose that for every node $C$, every phase-state $\xi\in\Xi_C$, and every full
action profile $a\in\prod_j A_j(\xi)$, the internal transition kernel
$P_C^0(\cdot\mid\xi,a)$ belongs to an irreducible family whose directed support
graph is the same strongly connected graph on $\Xi_C$.  Then every convex mixture
of those kernels is irreducible as well, so condition~\textnormal{(a)} holds
trivially.
\end{example}

\subsection{A uniform irreducibility modulus}

For fixed $\eta$, define
\begin{equation}
\label{eq:kappa-eta}
\kappa_\eta
:=
\min_{C\in\cT_\eta}
\min_{i\in N}
\inf_{\sigma_{-i}\in\prod_{j\ne i}\Sigma_{j,C}}
\min_{a_i(\cdot)}
\min_{\xi,\zeta\in\Xi_C}
\max_{1\le n\le |\Xi_C|-1}
\Bigl(Q_{i,C}^{\sigma_{-i},a_i}\Bigr)^n(\zeta\mid\xi),
\end{equation}
where $Q_{i,C}^{\sigma_{-i},a_i}$ is the unilateral internal kernel from
condition~\textnormal{(a)}.

\begin{lemma}
\label{lem:kappa-positive}
Assume condition~\textnormal{(a)}.  Then $\kappa_\eta>0$ for every $\eta\in(0,1)$.
\end{lemma}

\begin{proof}
Fix $C$, $i$, and a pure action selection $a_i(\cdot)$.  By
condition~\textnormal{(a)}, every kernel
$Q_{i,C}^{\sigma_{-i},a_i}$ is irreducible on the finite set $\Xi_C$.  Therefore,
for every pair of states $(\xi,\zeta)$ there exists a path of length at most
$|\Xi_C|-1$ from $\xi$ to $\zeta$ with positive probability.  Hence the displayed
quantity in \eqref{eq:kappa-eta} is strictly positive for each fixed
$\sigma_{-i}$.

The map $\sigma_{-i}\mapsto Q_{i,C}^{\sigma_{-i},a_i}$ is continuous by
\Cref{lem:coeff-affine}, and the compact domain
$\prod_{j\ne i}\Sigma_{j,C}$ is a finite product of simplices.  Therefore the
minimum over that domain is attained and remains strictly positive.  Finally there
are only finitely many nodes, players, and pure action selections, so the overall
minimum is strictly positive.
\end{proof}

\subsection{Bias-span bound}

\begin{lemma}[Expected hitting-time bound]
\label{lem:hitting-time}
Let $Q$ be an irreducible Markov kernel on a finite set $X$, and define
\[
\kappa(Q):=\min_{x,y\in X}\max_{1\le n\le |X|-1}Q^n(y\mid x).
\]
Then for all $x,y\in X$,
\[
\E_x[\tau_y]\le \frac{|X|-1}{\kappa(Q)},
\qquad
\tau_y:=\inf\{t\ge 0:X_t=y\}.
\]
\end{lemma}

\begin{proof}
Set $m:=|X|-1$.  By definition of $\kappa(Q)$,
\[
\Pb_x(\tau_y\le m)\ge \kappa(Q).
\]
Applying the Markov property after each block of $m$ steps gives
\[
\Pb_x(\tau_y>km)\le (1-\kappa(Q))^k
\qquad (k\ge 0).
\]
Therefore
\[
\E_x[\tau_y]
=
\sum_{t\ge 0}\Pb_x(\tau_y>t)
\le
m\sum_{k\ge 0}\Pb_x(\tau_y>km)
\le
m\sum_{k\ge 0}(1-\kappa(Q))^k
=
\frac{m}{\kappa(Q)}.
\]
\end{proof}

\begin{proposition}[Uniform bias-span bound]
\label{prop:bias-span}
Assume condition~\textnormal{(a)}.  Let $x^*\in\BR_\eta(x^*)$ be a fixed point and
let $(g_{i,C}^*,h_{i,C}^*)$ be the Bellman certificate of player $i$ in node $C$,
normalized by $\min_\xi h_{i,C}^*(\xi)=0$.  Then
\[
\spn(h_{i,C}^*)\le H_\eta:=\frac{\max_D|\Xi_D|-1}{\kappa_\eta}.
\]
\end{proposition}

\begin{proof}
Fix $C$ and $i$.  Let $\zeta\in\Xi_C$ attain the minimum of $h_{i,C}^*$, so
$h_{i,C}^*(\zeta)=0$.  Let $a_i^*(\cdot)$ be a measurable selector choosing, in
each state, an action where equality holds in the Bellman system.  The corresponding
unilateral internal kernel is irreducible by condition~\textnormal{(a)}.

For the Markov chain governed by that selector, the Bellman equality gives
\[
\E\bigl[h_{i,C}^*(X_{t+1})-h_{i,C}^*(X_t)\mid \cF_t\bigr]
\le g_{i,C}^*-u_i^t
\le 1,
\]
because $u_i^t\in[0,1]$ and $g_{i,C}^*\in[0,1]$.  Stopping at $\tau_\zeta$ and
summing yields
\[
h_{i,C}^*(\xi)-h_{i,C}^*(\zeta)\le \E_\xi[\tau_\zeta].
\]
By \Cref{lem:hitting-time} and \Cref{lem:kappa-positive},
\[
\E_\xi[\tau_\zeta]\le \frac{|\Xi_C|-1}{\kappa_\eta}
\le \frac{\max_D|\Xi_D|-1}{\kappa_\eta}=H_\eta.
\]
Since the minimum of $h_{i,C}^*$ is $0$, this proves the claimed span bound.
\end{proof}

\subsection{Deviation-value Lipschitz bound}

\begin{lemma}[Coefficient-level Bellman transfer]
\label{lem:coefficient-transfer}
Let $r,Q$ and $r',Q'$ be two finite average-reward MDP coefficient systems on the
same state-action graph, and let $(g,h)$ satisfy
\[
g+h(x)\ge r(x,a)+\sum_y Q(y\mid x,a)h(y)
\qquad\text{for all }(x,a).
\]
Then
\[
g+h(x)
\ge
r'(x,a)+\sum_y Q'(y\mid x,a)h(y)
-\Bigl(\|r'-r\|_\infty + \spn(h)\sup_{x,a}\|Q'(\cdot\mid x,a)-Q(\cdot\mid x,a)\|_{\mathrm{TV}}\Bigr)
\]
for all $(x,a)$.  In particular the perturbation constant is governed by
$1+\spn(h)$, not merely by $\spn(h)$.
\end{lemma}

\begin{proof}
Fix $(x,a)$.  Write
\begin{align*}
&r'(x,a)+\sum_y Q'(y\mid x,a)h(y)
-
\Bigl(r(x,a)+\sum_y Q(y\mid x,a)h(y)\Bigr)
\\
&= \bigl(r'(x,a)-r(x,a)\bigr)
+
\sum_y\bigl(Q'(y\mid x,a)-Q(y\mid x,a)\bigr)h(y).
\end{align*}
The first term is bounded by $\|r'-r\|_\infty$.  For the second term, subtract the
minimum of $h$ and use that signed measures of total mass $0$ have dual norm equal
to total variation against functions of oscillation $\spn(h)$:
\[
\left|\sum_y\bigl(Q'-Q\bigr)(y\mid x,a)h(y)\right|
\le
\spn(h)\,\|Q'(\cdot\mid x,a)-Q(\cdot\mid x,a)\|_{\mathrm{TV}}.
\]
Combining the two bounds gives the claim.
\end{proof}

\begin{lemma}[Deviation-value Lipschitz bound]
\label{lem:deviation-lipschitz}
Fix $\eta\in(0,1)$ and a node $C$.  There exists a finite constant $L_{\eta,C}$
such that for every player $i$, every two opponents' kernels
$\sigma_{-i},\tau_{-i}\in\prod_{j\ne i}\Sigma_{j,C}$, and every continuation vector
$v\in[0,1]^{B_C\times N}$,
\[
\bigl|\operatorname{Val}_{i,C}(\tau_{-i},v)-\operatorname{Val}_{i,C}(\sigma_{-i},v)\bigr|
\le L_{\eta,C}\,\|\tau_{-i}-\sigma_{-i}\|_1,
\]
where $\operatorname{Val}_{i,C}(\cdot,v)$ denotes the optimal average reward of the
deviation MDP.  One may take
\[
L_{\eta,C}=K_{i,C}(1+H_\eta),
\]
with $K_{i,C}$ as in \Cref{lem:coeff-affine}.
\end{lemma}

\begin{proof}
Fix $\sigma_{-i}$ and let $(g,h)$ be an optimal Bellman certificate for the MDP with
coefficients induced by $(\sigma_{-i},v)$.  By \Cref{prop:bias-span},
$\spn(h)\le H_\eta$.  Apply \Cref{lem:coefficient-transfer} with the primed
coefficients induced by $(\tau_{-i},v)$.  Using the Lipschitz estimates from
\Cref{lem:coeff-affine} yields
\[
\operatorname{Val}_{i,C}(\tau_{-i},v)
\le
\operatorname{Val}_{i,C}(\sigma_{-i},v)
+
K_{i,C}(1+H_\eta)\,\|\tau_{-i}-\sigma_{-i}\|_1.
\]
Interchanging the roles of $\sigma_{-i}$ and $\tau_{-i}$ gives the reverse
inequality.
\end{proof}

\subsection{Blockwise Bellman inequality}

Fix a fixed point $x^*\in\BR_\eta(x^*)$ and the realized profile
$\sigma_{\eta,\delta,M_*}$ from \Cref{thm:realization}.  The key point is that in a
complete block played in node $C$, the compliant profile is close to the local
stationary kernel $\sigma_C^*$, so the local Bellman certificate transfers to the
realized block.

\begin{lemma}
\label{lem:block-bellman}
There exists a constant $C_\eta<\infty$ such that for every complete block $B$ of
length $|B|\ge M_*$ played in node $C$, every player $i$, and every unilateral
deviation $\tau_i$,
\begin{equation}
\label{eq:block-bellman-ineq}
\E\left[
\frac1{|B|}\sum_{t\in B}u_i^t(\tau_i,\sigma_{-i})
\Bigm|\cF_{\mathrm{in}(B)}
\right]
\le
 g_{i,C}^* + \eta + C_\eta\delta + \frac{2H_\eta}{|B|}.
\end{equation}
\end{lemma}

\begin{proof}
Fix $B$, $C$, and $i$.  The Bellman certificate at the fixed point gives, for every
phase-state $\xi$ and every action $a_i\in A_i(\xi)$,
\begin{align}
\label{eq:exact-bellman-fixed-point}
g_{i,C}^* + h_{i,C}^*(\xi)
\ge {}
&r_{i,C}^{\sigma_{-i,C}^*,x^*}(\xi,a_i)
+
\sum_{\xi'\in\Xi_C}Q_{i,C}^{\sigma_{-i,C}^*}(\xi'\mid\xi,a_i)h_{i,C}^*(\xi')
\\
&+
\sum_{b\in B_C}\Lambda_{i,C}^{\sigma_{-i,C}^*}(b\mid\xi,a_i)v_{i,C}^{x^*}(b)
-\eta.
\nonumber
\end{align}

During the realized block, the compliant controller is $\delta$-close in expected
augmented occupation to the fixed-point node controller by
\eqref{eq:block-occupation-realization}.  Since the coefficient arrays of the
deviation MDP are Lipschitz in the compliant kernel, \Cref{lem:deviation-lipschitz}
shows that the one-step coefficients seen by the deviator differ from the fixed
coefficients in \eqref{eq:exact-bellman-fixed-point} by at most $C_\eta\delta$ in
expectation, for a uniform constant $C_\eta$ depending only on the finite tree
object and $H_\eta$.  Therefore the process
\[
M_t:=h_{i,C}^*(X_t)-\sum_{s=t_0}^{t-1}
\bigl(u_i^s(\tau_i,\sigma_{-i})-g_{i,C}^* -\eta - C_\eta\delta\bigr)
\]
has nonnegative conditional drift along the block.  Summing from the block entry
$t_0$ to the block exit $t_1-1$ and taking expectations yields
\[
\E\left[\sum_{t\in B}u_i^t(\tau_i,\sigma_{-i})\Bigm|\cF_{\mathrm{in}(B)}\right]
\le
|B|\bigl(g_{i,C}^*+\eta+C_\eta\delta\bigr)
+
\E\bigl[h_{i,C}^*(X_{t_0})-h_{i,C}^*(X_{t_1})\mid\cF_{\mathrm{in}(B)}\bigr].
\]
The terminal bias difference is bounded in absolute value by
$2\spn(h_{i,C}^*)\le 2H_\eta$ by \Cref{prop:bias-span}.  Dividing by $|B|$ gives
\eqref{eq:block-bellman-ineq}.
\end{proof}

\begin{corollary}
\label{cor:block-regret}
There exists a constant $C_\eta<\infty$ such that for every complete block $B$ of
length at least $M_*$ played in node $C$, every player $i$, and every unilateral
deviation $\tau_i$,
\[
\E\left[
\frac1{|B|}\sum_{t\in B}u_i^t(\tau_i,\sigma_{-i})
-
\frac1{|B|}\sum_{t\in B}u_i^t(\sigma)
\Bigm|\cF_{\mathrm{in}(B)}
\right]
\le
\eta + C_\eta\delta + \frac{2H_\eta}{|B|}.
\]
\end{corollary}

\begin{proof}
Subtract the realized compliant block average from
\eqref{eq:block-bellman-ineq} and use
\eqref{eq:block-payoff-realization}.
\end{proof}

\subsection{Proof of the main theorem}

\begin{proof}[Proof of \Cref{thm:main}]
Fix $\eps>0$.  Choose parameters in the order
\[
\eta,\qquad \delta,\qquad M_*,\qquad T_0.
\]

\smallskip
\noindent\textbf{Step 1: fixed point.}
Choose $\eta>0$ with $\eta\le \eps/4$.  By
\Cref{thm:kakutani-global}, there exists a fixed point $x^*\in\BR_\eta(x^*)$.

\smallskip
\noindent\textbf{Step 2: bias span.}
By \Cref{lem:kappa-positive}, $\kappa_\eta>0$.  Define
\[
H_\eta:=\frac{\max_D|\Xi_D|-1}{\kappa_\eta}.
\]
By \Cref{prop:bias-span}, every fixed-point Bellman bias has span at most $H_\eta$.

\smallskip
\noindent\textbf{Step 3: realization accuracy.}
Choose $\delta>0$ such that $C_\eta\delta\le \eps/4$, where $C_\eta$ is the
constant from \Cref{cor:block-regret}.  Freeze a superexponential schedule at a
level $R$ so large that the associated tail block length $M_*=M_R$ satisfies
\[
\frac{2H_\eta}{M_*}\le \frac{\eps}{4}.
\]
Let $\sigma^\eps:=\sigma_{\eta,\delta,M_*}$ be the realized profile from
\Cref{thm:realization}.

\smallskip
\noindent\textbf{Step 4: horizon threshold.}
Choose $T_0$ so large that
\[
\frac{K_\eta}{T_0}\le \frac{\eps}{4},
\]
with $K_\eta$ from \Cref{thm:realization}.

Now fix $T\ge T_0$, an initial state $s$, a player $i$, and a unilateral deviation
$\tau_i$.  Partition the horizon $\{1,\dots,T\}$ into the complete blocks
contained in $[1,T]$ plus the incomplete initial and terminal fragments.  Applying
\Cref{lem:aggregation} from Appendix~\ref{app:aggregation} with
\[
a=\eta+C_\eta\delta,
\qquad H=H_\eta,
\]
and using \Cref{cor:block-regret} yields
\[
\gamma_i^T\bigl(s,(\tau_i,\sigma^\eps_{-i})\bigr)
-
\gamma_i^T(s,\sigma^\eps)
\le
\eta + C_\eta\delta + \frac{2H_\eta}{M_*} + \frac{K_\eta}{T}.
\]
By construction, the right-hand side is at most
\[
\frac{\eps}{4}+\frac{\eps}{4}+\frac{\eps}{4}+\frac{\eps}{4}=\eps.
\]
Therefore
\[
\gamma_i^T(s,\sigma^\eps)
\ge
\gamma_i^T\bigl(s,(\tau_i,\sigma^\eps_{-i})\bigr)-\eps.
\]
Since $T\ge T_0$, $s$, $i$, and $\tau_i$ were arbitrary, $\sigma^\eps$ is a
public finite-memory uniform $\eps$-equilibrium.
\end{proof}

\section{Discussion}

The full Neyman conjecture remains open.  What the present paper does is narrower
but mathematically sharp: it shows that the downstream compact-tree route is
correct once the finite universal $\eta$-tree object is available and once the
nodewise irreducibility obstruction from condition~\textnormal{(a)} is excluded.

The most important conceptual point is the augmented-node repair.  The raw
stationarity equations from v4 were not a minor typo; they annihilated all exit
mass.  The regenerated augmented system restores the correct balance and makes the
graph theorem for fluxes precise.  The second important point is the distinction
between \emph{frozen} continuation values and \emph{self-consistent} continuation
values.  The response correspondence must be defined against frozen data; only the
fixed point enforces self-consistency.  Finally, the local controller estimate
never needed an exponential rate.  The polynomial $L^{-1/2}$ rate is enough, and it
fits the block-aggregation argument exactly.

\appendix

\section{Finite average-reward MDP facts}
\label{app:mdp}

This appendix records the one-player linear-program facts used in the main text.
All statements are standard; see Puterman's monograph \cite{Puterman1994} for
background.

Consider a finite average-reward MDP with state space $X$, action sets $A(x)$,
reward function $r(x,a)\in[0,1]$, and transition kernel $Q(\cdot\mid x,a)$.

\begin{proposition}[Primal-dual formulation]
\label{prop:mdp-duality}
The optimal average reward is the common value of the primal and dual linear
programs
\begin{align*}
\textnormal{maximize }&\sum_{x\in X}\sum_{a\in A(x)}\nu(x,a)r(x,a)
\\
\textnormal{subject to }&\nu\ge 0,
\qquad \sum_{x,a}\nu(x,a)=1,
\\
&\sum_{a\in A(y)}\nu(y,a)=\sum_{x,a}\nu(x,a)Q(y\mid x,a)
\qquad (y\in X),
\end{align*}
and
\begin{align*}
\textnormal{minimize }&g
\\
\textnormal{subject to }&g+h(x)\ge r(x,a)+\sum_y Q(y\mid x,a)h(y)
\qquad (x\in X,\ a\in A(x)).
\end{align*}
There exists an optimal stationary mixed policy and an optimal dual pair $(g,h)$.
Moreover, after normalization $\min_x h(x)=0$, one may choose $h$ in a compact box.
\end{proposition}

\begin{proof}
This is the classical average-reward LP formulation for finite MDPs.  Feasibility is
immediate.  The primal feasible set is a nonempty compact polytope, so an optimum
exists.  The dual feasible set is nonempty because choosing $g=1$ and $h\equiv 0$
is feasible.  Finite-dimensional LP duality gives equality of primal and dual
optima and existence of optimal primal and dual solutions.  The occupation measure
of an optimal primal solution induces an optimal stationary mixed policy.  The
normalization $\min_x h(x)=0$ is obtained by subtracting a constant.
\end{proof}

\begin{remark}
The proof of \Cref{prop:local-BR} uses this proposition with the local deviation
coefficients from Section~\ref{sec:local}.  No multiplayer statement is proved in
this appendix; the multiplayer part enters only through Kakutani in
\Cref{prop:local-equilibrium} and \Cref{thm:kakutani-global}.
\end{remark}

\section{A uniform block estimate for finite Markov chains}
\label{app:block-estimate}

The local controller lemma uses only a polynomial empirical-average bound.

\begin{proposition}
\label{prop:block-estimate}
Let $Q$ be an \emph{irreducible} Markov kernel on a finite state space $X$, and let
$\pi$ be its unique invariant distribution.  There exists a constant $K(Q)<\infty$
such that for every bounded function $f:X\to[-1,1]$, every initial state $x\in X$,
and every $L\ge 1$,
\[
\E_x\left[\left|\frac1L\sum_{t=1}^L f(X_t)-\sum_{y\in X}\pi(y)f(y)\right|\right]
\le \frac{K(Q)}{\sqrt L}.
\]
\end{proposition}

\begin{proof}
For an irreducible finite Markov chain the invariant distribution is unique and the
Poisson equation
\[
\psi-Q\psi = f-\pi(f)
\]
has a bounded solution $\psi$ after normalizing by $\pi(\psi)=0$.  Then
\[
\sum_{t=1}^L \bigl(f(X_t)-\pi(f)\bigr)
=
\psi(X_1)-\psi(X_{L+1}) + M_L,
\]
where $(M_L)$ is a martingale with bounded increments depending only on $Q$ and on
$\|f\|_\infty\le 1$.  Taking expectations, using the boundedness of $\psi$, and
applying the Cauchy--Schwarz inequality to the martingale term gives an
$L^{-1/2}$ bound with a constant depending only on $Q$.
\end{proof}

\section{Aggregating block estimates}
\label{app:aggregation}

\begin{lemma}
\label{lem:aggregation}
Let the horizon $\{1,\dots,T\}$ be partitioned into pairwise disjoint intervals
\[
I^{\mathrm{pre}},\ B_1,\dots,B_m,\ I^{\mathrm{post}},
\]
where $B_1,\dots,B_m$ are the complete blocks and the other two intervals are the
possibly empty incomplete fragments.  Assume there exist constants $a\ge 0$, $H\ge
0$, $M_*\ge 1$, and $K<\infty$ such that:
\begin{enumerate}[label=\textnormal{(\roman*)},leftmargin=2.6em]
\item every complete block satisfies $|B_r|\ge M_*$;
\item for every $r$,
\[
\E\left[\sum_{t\in B_r}\Delta_t\Bigm|\cF_{\mathrm{in}(B_r)}\right]
\le a|B_r|+2H,
\]
where $\Delta_t$ is the stagewise regret increment;
\item the incomplete fragments satisfy
\[
\E\left[\sum_{t\in I^{\mathrm{pre}}\cup I^{\mathrm{post}}}\Delta_t\right]\le K.
\]
\end{enumerate}
Then
\[
\frac1T\E\left[\sum_{t=1}^T\Delta_t\right]
\le a + \frac{2H}{M_*}+\frac{K}{T}.
\]
\end{lemma}

\begin{proof}
Take expectations in the blockwise estimate and sum over the complete blocks:
\[
\E\left[\sum_{r=1}^m\sum_{t\in B_r}\Delta_t\right]
\le a\sum_{r=1}^m|B_r| + 2Hm.
\]
Since each complete block has length at least $M_*$,
\[
m\le \frac{1}{M_*}\sum_{r=1}^m |B_r|.
\]
Hence
\[
\E\left[\sum_{r=1}^m\sum_{t\in B_r}\Delta_t\right]
\le
\left(a+\frac{2H}{M_*}\right)\sum_{r=1}^m|B_r|.
\]
Adding the incomplete fragments and using assumption~(iii) gives
\[
\E\left[\sum_{t=1}^T\Delta_t\right]
\le
\left(a+\frac{2H}{M_*}\right)T + K.
\]
Divide by $T$.
\end{proof}

\begin{thebibliography}{99}

\bibitem{BewleyKohlberg1976}
T.~Bewley and E.~Kohlberg.
\newblock The asymptotic theory of stochastic games.
\newblock \emph{Mathematics of Operations Research}, 1(3):197--208, 1976.

\bibitem{FleschThuijsmanVrieze1997}
J.~Flesch, F.~Thuijsman, and K.~Vrieze.
\newblock Cyclic {M}arkov equilibria in stochastic games.
\newblock \emph{International Journal of Game Theory}, 26:303--314, 1997.

\bibitem{Kakutani1941}
S.~Kakutani.
\newblock A generalization of {B}rouwer's fixed point theorem.
\newblock \emph{Duke Mathematical Journal}, 8(3):457--459, 1941.

\bibitem{MertensNeyman1981}
J.-F.~Mertens and A.~Neyman.
\newblock Stochastic games.
\newblock \emph{International Journal of Game Theory}, 10(2):53--66, 1981.

\bibitem{MertensSorinZamir2003}
J.-F.~Mertens, S.~Sorin, and S.~Zamir.
\newblock Stochastic games.
\newblock In R.~Aumann and S.~Hart, editors, \emph{Handbook of Game Theory with
Economic Applications}, volume~3, chapter~47, pages 1811--1832. Elsevier,
Amsterdam, 2003.

\bibitem{Neyman2003}
A.~Neyman.
\newblock From {M}arkov chains to stochastic games.
\newblock In A.~Neyman and S.~Sorin, editors, \emph{Stochastic Games and
Applications}, NATO Science Series C, volume~570, pages 9--25. Kluwer Academic
Publishers, Dordrecht, 2003.

\bibitem{Nowak2003}
A.~S.~Nowak.
\newblock N-person stochastic games: Extensions of the finite state space case and
correlation.
\newblock In A.~Neyman and S.~Sorin, editors, \emph{Stochastic Games and
Applications}, NATO Science Series C, volume~570, pages 93--106. Kluwer Academic
Publishers, Dordrecht, 2003.

\bibitem{Puterman1994}
M.~L.~Puterman.
\newblock \emph{Markov Decision Processes: Discrete Stochastic Dynamic
Programming}.
\newblock John Wiley \& Sons, New York, 1994.

\bibitem{Shapley1953}
L.~S.~Shapley.
\newblock Stochastic games.
\newblock \emph{Proceedings of the National Academy of Sciences of the United
States of America}, 39(10):1095--1100, 1953.

\bibitem{Solan1999}
E.~Solan.
\newblock Three-player absorbing games.
\newblock \emph{Mathematics of Operations Research}, 24(3):669--698, 1999.

\bibitem{Solan2003}
E.~Solan.
\newblock Uniform equilibrium: More than two players.
\newblock In A.~Neyman and S.~Sorin, editors, \emph{Stochastic Games and
Applications}, NATO Science Series C, volume~570, pages 285--294. Kluwer Academic
Publishers, Dordrecht, 2003.

\bibitem{SolanVieille2002}
E.~Solan and N.~Vieille.
\newblock Correlated equilibrium in stochastic games.
\newblock \emph{Games and Economic Behavior}, 38(2):362--399, 2002.

\bibitem{Vieille2000I}
N.~Vieille.
\newblock Equilibrium in two-person stochastic games {I}: A reduction.
\newblock \emph{Israel Journal of Mathematics}, 119:55--91, 2000.

\bibitem{Vieille2000II}
N.~Vieille.
\newblock Equilibrium in two-person stochastic games {II}: The case of recursive
games.
\newblock \emph{Israel Journal of Mathematics}, 119:93--126, 2000.

\bibitem{Vieille2002}
N.~Vieille.
\newblock Stochastic games: Recent results.
\newblock In R.~Aumann and S.~Hart, editors, \emph{Handbook of Game Theory with
Economic Applications}, volume~3, chapter~48, pages 1833--1850. Elsevier,
Amsterdam, 2002.

\end{thebibliography}

\end{document}
